%--------------------------------------------------------------------------------
%
%
%--------------------------------------------------------------------------------
\pagebreak
\unnumberedSection{FxCA - flush from Instruction or Data Cache}

\begin{iPageDescEntry}{Syntax}
	\texttt{FICA [ RegX ] ( RegB )} \\
	\texttt{FDCA [ RegX ] ( RegB )}
\end{iPageDescEntry}

\begin{iPageDescEntry}{Format}
 	The \texttt{FICA} instruction uses the following instruction format. \\
    \begin{flushleft}
    \drawInstrFormat
        {0,1,3,5,6.5,8.5,9.5,11.5,16}
        {\OpCodeGrpSYS, \OpCodePCA, 0, 0, RegB, 0, RegX, 0}
    \end{flushleft}
    The \texttt{FDCA} instruction uses the following instruction format. \\
    \begin{flushleft}
    \drawInstrFormat
        {0,1,3,5,6.5,8.5,9.5,11.5,16}
        {\OpCodeGrpSYS, \OpCodePCA, 0, 1, RegB, 0, RegX, 0}
    \end{flushleft}
\end{iPageDescEntry}

\begin{iPageDescEntry}{Description}        
    The \texttt{FICA} instruction invalidates a cache line, selected by the effective address formed by adding \texttt{RegX} to \texttt{RegB}. 
    
    The \texttt{FDCA} instruction flushes and invalidates a cache line, selected by the effective address formed by adding \texttt{RegX} to \texttt{RegB}. The cache line addressed is written back to memory if and only if it is dirty, and then invalidated. 

    Both instructions are privileged instructions.
\end{iPageDescEntry}

\begin{iPageDescOpEntry}{Operation}
\begin{lstlisting}[style=iPageOpStyle]
flushFromCache( addOfs32( RegB, RegX ));
\end{lstlisting}
\end{iPageDescOpEntry}

\begin{iPageDescEntry}{Exceptions}
    Privilege violation trap.
\end{iPageDescEntry}

\begin{iPageDescEntry}{Notes}
None.
\end{iPageDescEntry} 
